\section{Conclusions}
\label{sec:conclusion}

\todolist{tim}{
\item constraints are useful to write reactive systems
\item the relationships between objects are modularized
\item adding new codepaths that affect the relationships works
  automatically}
  
%% softened the conclusion a little - I don't think we have shown this
%% conclusively yet
We have argued that \ac{ocp} facilitates reactive systems, because it
modularizes the relationship between objects and decouples constraint
satisfaction from the application.  In Babelsberg, constraints can be
dynamically added and removed and maintained automatically, making
them useful for writing interactive applications.  As an example, we
implemented an interactive application for specifying and solving a
Sudoku with constraints.  With the interactive \ac{gui}, the user can
change the values of the constraint variables without breaking the
properties of the Sudoku.  The application reacts on the user input by
resolving the underlying constraints.
%%% Local Variables:
%%% mode: latex
%%% TeX-master: "paper"
%%% End:
